\section{Introduction}
\label{sec:introduction}

Denoising diffusion models generate images by simulating a reverse trajectory
from noise to data. This simulation is accurate when many denoiser evaluations
are available, but those evaluations dominate sampling cost. Fast sampling
therefore asks a concrete allocation question: given only a small number of
function evaluations, where along the reverse trajectory should they be placed?
The answer is not determined by the solver alone. Even a strong numerical
update can fail if its few steps are spent in the wrong regions of the
denoising path.

Discretization schedules are the mechanism for making this allocation. Linear,
polynomial, log-SNR, and EDM-style schedules are simple, stable, and widely
used \citep{karras2022edm}, but they are usually hand-crafted and reused across
models, solvers, and budgets. Recent work shows that this choice can be
optimized. AYS searches for schedules with a solver-specific Girsanov KL upper
bound \citep{sabour2024ays}. Score-Optimal Diffusion Schedules derive a local
path cost from the work performed by a score-based corrector and choose an
inverse path-length schedule \citep{williams2024scoreoptimal}. Other
optimized-timestep methods minimize discrepancies between ODE teacher
solutions and numerical solver trajectories \citep{xue2024optimizedtimesteps}.
These approaches make clear that schedule selection is a central component of
fast diffusion sampling.

However, the criteria used to choose a schedule often remain indirect from the
perspective of the sampler that will actually be run. A distributional bound,
transport-style cost, or path information increment can justify an effective
grid. These criteria leave open the diagnostic question we care about: where
does this numerical solver make its largest local prediction errors along the
denoising trajectory? Without that information, schedule selection still looks
like global search instead of targeted correction of the solver's failure modes.

We present \emph{geometric prediction defect equalization}, a training-free
algorithm for automatically selecting schedules for deterministic diffusion ODE
samplers. Our method uses accurate teacher trajectories from a pretrained model
to estimate where a target student solver locally deviates
from the intended trajectory. It then allocates time steps so that these local
geometric prediction defects are equalized across the schedule.

The details of our method are designed to make this diagnostic meaningful. Each
defect probe starts from the correct teacher state, so it measures one-step
solver error rather than error accumulated from previous approximate steps.
Teacher trajectories are treated as Monte Carlo samples from the teacher
marginal at each noise level, so the estimated defect profile is
distributional rather than path-specific. A state-space metric and optional
tangent downweighting turn endpoint residuals into geometric prediction
defects. Finally, a local power-law model determines how to convert these
defects into a monitor density whose inverse CDF yields the schedule.

Our contributions are as follows.
\begin{enumerate}
  \item \textbf{A solver-facing schedule objective.} We formulate low-NFE
  schedule selection as distributional local prediction-defect equalization for
  a fixed pretrained model and deterministic sampler.
  \item \textbf{A training-free calibration algorithm.} We estimate local
  oracle-start defects from teacher trajectories, aggregate them over the
  teacher marginal with mean, CVaR, or mixed risk, and construct the schedule
  by inverse-CDF equalization of the induced monitor.
  \item \textbf{Geometry and theory.} We derive the monitor
  $\omega_{\mathcal R}(u)=(\bar a_{\mathcal R}(u)+\epsilon_a)^{1/q}$ from a
  local power-law defect model, prove coordinate covariance of the monitor
  one-form, and prove asymptotic minimax optimality of the inverse-CDF schedule
  for local distributional risk.
  \item \textbf{Scoped empirical evidence.} We report a real PNDM/CIFAR-10
  smoke gate, three-seed 50k CIFAR-10 FID results including base, linear,
  Karras-style fixed, and authorized project-owned offline-proxy schedules,
  matched SD1.5 and SDXL CFG sweeps retained as text-to-image boundary
  evidence, and model-evaluation-equivalent calibration cost accounting.
\end{enumerate}

The empirical claims are deliberately bounded. On PNDM/CIFAR-10 with the Euler
solver, our schedule improves FID over base, native-coordinate linear, and the
project-owned offline-proxy schedule at NFE 10 and 20. The comparison to a
fixed Karras/EDM-style grid is mixed: ours wins at NFE 10 and trails
slightly at NFE 20. The matched text-to-image sweeps show condition-dependent
behavior: SD1.5 is negative under these metrics, while SDXL improves
ImageReward against base at CFG 5.0 and 7.5 and has positive mean deltas
against AYS at CFG 7.5. We therefore present our method as a way to expose and
equalize solver-local
trajectory defects; direct optimization of FID, CLIPScore, ImageReward, or
human preference remains outside the claim.
